\documentclass{resume}

\newcommand{\en}[1]{}
\newcommand{\zh}[1]{#1}
\geometry{left=0.45in,right=0.45in,top=0.14in,bottom=0.20in}
\setlist[itemize]{itemsep=0.15ex,topsep=0.15em,leftmargin=1.25pc}

\zh{
  \usepackage{xeCJK}
  % Helper to set main/sans/mono CJK fonts consistently without triplicating options.
  \newcommand{\resumeCJKFaces}[2]{%
    \setCJKmainfont{#1}[#2]%
    \setCJKsansfont{#1}[#2]%
    \setCJKmonofont{#1}[#2]%
  }
  % SimSun has no real bold/italic; FakeBold on upright matches Medium body weight,
  % while stronger FakeBold on bold keeps \textbf distinct.
  \newcommand{\resumeCJKFakeBoldFaces}[1]{%
    \resumeCJKFaces{#1}{%
      UprightFont=*, BoldFont=*, ItalicFont=*, BoldItalicFont=*,
      UprightFeatures={FakeBold=1.8},
      BoldFeatures={FakeBold=4.6},
      ItalicFeatures={FakeBold=1.8},
      BoldItalicFeatures={FakeBold=4.6}}%
  }
  \IfFontExistsTF{SimSun}{
    \resumeCJKFakeBoldFaces{SimSun}
  }{
    \IfFontExistsTF{Songti SC}{
      \resumeCJKFaces{Songti SC}{%
        UprightFont=* Regular, BoldFont=* Bold,
        ItalicFont=* Regular, BoldItalicFont=* Bold,
        UprightFeatures={FakeBold=1.2},
        ItalicFeatures={FakeBold=1.2}}
    }{
      \IfFontExistsTF{Noto Serif CJK SC Medium}{
        \resumeCJKFaces{Noto Serif CJK SC}{%
          UprightFont=* Medium, BoldFont=* Bold,
          ItalicFont=* Medium, BoldItalicFont=* Bold}
      }{
        \IfFontExistsTF{Noto Serif CJK SC}{
          \resumeCJKFaces{Noto Serif CJK SC}{%
            UprightFont=* Regular, BoldFont=* Bold,
            ItalicFont=* Regular, BoldItalicFont=* Bold,
            UprightFeatures={FakeBold=1.2},
            ItalicFeatures={FakeBold=1.2}}
        }{
          \resumeCJKFaces{FandolSong-Regular}{%
            Extension=.otf,
            UprightFont=FandolSong-Regular, BoldFont=FandolSong-Bold,
            ItalicFont=FandolSong-Regular, BoldItalicFont=FandolSong-Bold,
            UprightFeatures={FakeBold=1.2},
            ItalicFeatures={FakeBold=1.2}}
        }
      }
    }
  }
  % Load TeX Live's bundled Latin Modern OpenType fonts by filename (kpathsea)
  % so macOS CoreText doesn't fail on OS font-name lookup, matching tulmr/tulmtt.
  \setmainfont{lmroman10-regular}[
    Extension = .otf,
    UprightFont = lmroman10-regular,
    BoldFont = lmroman10-bold,
    ItalicFont = lmroman10-italic,
    BoldItalicFont = lmroman10-bolditalic,
    SmallCapsFont = lmromancaps10-regular,
    UprightFeatures = {
      FakeBold = 1.4,
      SizeFeatures = {
        {Size={-5.5},    Font=lmroman5-regular},
        {Size={5.5-6.5}, Font=lmroman6-regular},
        {Size={6.5-7.5}, Font=lmroman7-regular},
        {Size={7.5-8.5}, Font=lmroman8-regular},
        {Size={8.5-9.5}, Font=lmroman9-regular},
        {Size={9.5-11},  Font=lmroman10-regular},
        {Size={11-15},   Font=lmroman12-regular},
        {Size={15-},     Font=lmroman17-regular}
      }
    },
    BoldFeatures = {
      SizeFeatures = {
        {Size={-5.5},    Font=lmroman5-bold},
        {Size={5.5-6.5}, Font=lmroman6-bold},
        {Size={6.5-7.5}, Font=lmroman7-bold},
        {Size={7.5-8.5}, Font=lmroman8-bold},
        {Size={8.5-9.5}, Font=lmroman9-bold},
        {Size={9.5-11},  Font=lmroman10-bold},
        {Size={11-},     Font=lmroman12-bold}
      }
    },
    ItalicFeatures = {
      FakeBold = 1.4,
      SizeFeatures = {
        {Size={-7.5},    Font=lmroman7-italic},
        {Size={7.5-8.5}, Font=lmroman8-italic},
        {Size={8.5-9.5}, Font=lmroman9-italic},
        {Size={9.5-11},  Font=lmroman10-italic},
        {Size={11-},     Font=lmroman12-italic}
      }
    },
    SmallCapsFeatures = {FakeBold = 1.4}]
  \setmonofont{lmmono10-regular}[
    Extension = .otf,
    UprightFont = lmmono10-regular,
    BoldFont = lmmonolt10-bold,
    ItalicFont = lmmono10-italic,
    BoldItalicFont = lmmonolt10-boldoblique,
    UprightFeatures = {
      FakeBold = 1.4,
      SizeFeatures = {
        {Size={-8.5},    Font=lmmono8-regular},
        {Size={8.5-9.5}, Font=lmmono9-regular},
        {Size={9.5-11},  Font=lmmono10-regular},
        {Size={11-},     Font=lmmono12-regular}
      }
    },
    ItalicFeatures = {FakeBold = 1.4}]
}

\begin{document}
\pagestyle{empty}
\fontsize{8.6pt}{13.0pt}\selectfont
\emergencystretch=1.5em
\vspace*{-0.8ex}

\noindent
\begin{minipage}[t]{0.74\textwidth}
  {\LARGE\scshape \en{Qingliu Wu}\zh{吴清柳}\par}
  \vspace{0.6ex}
  {\ttfamily\normalsize
    \phone{17353675760} \textperiodcentered\
    \email{wql@bupt.edu.cn} \textperiodcentered\
    \github[micuks]{https://github.com/micuks}\par
  }
\end{minipage}\hfill
\makebox[0pt][r]{\raisebox{-7ex}[0pt][0pt]{\includegraphics[width=2.2cm,clip,trim=0 57 0 0]{selfie.png}}}
\vspace{-0.2ex}

\section{\en{Education}\zh{教育经历}}
\en{\datedsubsection{\textbf{Beijing University of Posts and Telecommunications}, Master Student}{09/2023 -- Present}}
\zh{\datedsubsection{\textbf{北京邮电大学}, 硕士研究生}{2024/09 -- 2027/06}}
\begin{itemize}[parsep=0.25ex]
	\item \zh{计算机科学与技术（学硕），获\textbf{2024--2025 年度国家奖学金}；IDMG 实验室，师从邵蓥侠老师，方向 \textbf{AI4DB}}
\end{itemize}
\en{\datedsubsection{\textbf{Beijing University of Posts and Telecommunications}, Bachelor of Engineering}{09/2020 -- 06/2024}}
\zh{\datedsubsection{\textbf{北京邮电大学}, 工学学士}{2020/09 -- 2024/06}}
\begin{itemize}[parsep=0.25ex]
	\item \zh{计算机科学与技术（计算机学院实验班），综合成绩 91.91，排名 29/422（前 6.56\%）}
\end{itemize}

\section{\zh{实习经历}}

\zh{\role{Qwen3.8-Next / Qwen3.6 训练引擎优化 · FP8 / MegaMoE / Megatron}{2026.10 -- 至今}}
\zh{{\small\itshape 快手 · 算法引擎部\ 训练引擎组 | AI Infra}\par}
\begin{itemize}[parsep=0.25ex]
	\item \zh{\textbf{Megatron FP8 训练优化：}面向 Qwen3.8-Next / Qwen3.6（GDN+Attn+MoE+MTP）262K 长序列训练，按 Hopper FP8 TN 布局管理 \texttt{X/dY/W} 生命周期：非 MoE \texttt{W\_col} 仅在 \texttt{dgrad} 前由 \texttt{W\_row} 免重量化转置、用完即释放，并跳过 \texttt{post\_all\_gather} 无效转置；\texttt{O-proj} 前向一次量化生成 \texttt{x\_row/x\_col}，反向仅保存 FP8 \texttt{x\_col}，不再保存 BF16 输入；删除 MegaMoE 不读取的专家 \texttt{W\_col}，再省约 2 GiB。配合 RoPE 按 CP rank 本地切分与 Attn 轻量重算，GDN+MoE 改为整层单 checkpoint，重算量基本不变、单层激活 168.6 $\to$ 4.2 MiB（$-$97.5\%）；35B-A3B 32 卡 H200 同布局峰值显存 $-$26.0\%（69.7 $\to$ 51.6 GiB）。}
	\item \zh{\textbf{MegaMoE 全链路接入与优化：}将 CuTe DSL 融合核 MegaMoE 接入 Qwen3.6 并迁移至 Qwen3.8-Next：将 FP32 主权重、FP8 参数与 \texttt{main\_grad} 统一为 \texttt{interleave8} 布局，Adam 更新后原地量化；跨 DP 的 $128{\times}128$ 量化块先 MAX 归约统一 scale，专家逆序注册抵消 DDP 建桶倒序，前向以连续 3D 视图零拷贝接入，去除每步 160+ 次拼接重量化（链路累计步时 $-$555 ms、显存 $-$352 MiB）；以 \texttt{as\_strided} 将连续等距的 DDP 梯度桶构造为 3D 视图，\texttt{wgrad} epilogue 跨 microbatch 直接累加 \texttt{main\_grad}，去除 128 个 cast/add kernel 与 384 MiB 临时梯度；合并前反向工作区，以方向私有 4B 备份槽保存/恢复 NVLink barrier 相位避免死锁，配合 Triton 紧凑路由、GPU 溢出检查与分块转置（48 $\to$ 6.5 ms），这部分合计步时约 $-$0.34 s、单卡峰值显存 $-$2.5 GiB。}
	\item \zh{\textbf{长序列算子融合与自研内核：}为 248K 大词表实现单遍 cross-entropy，梯度写回 \texttt{main\_grad}（1.62$\times$），并做 SFT 有效行压缩与 CP 均衡；融合 QK-Norm 与 partial mRoPE（3.09$\times$），优化 GDN 重算复用、\texttt{tile16} 与 AllToAll 融合；接入自研 SM90 两阶段 TMA 注意力反向核（较 FA3 快 15\%--21\%，\texttt{H=1} 时省 1.1 GiB），32 卡 262K 23 步精度对齐，同布局步时 $-$15.9\%（23.43 $\to$ 19.72 s）。}
\end{itemize}

\zh{\role{多模态大模型 Keye 预训练 · Megatron/DeepSpeed 框架优化}{2026.05 -- 2026.10}}
\zh{{\small\itshape 快手 · 可灵 AI 事业部 | AI Infra}\par}
\begin{itemize}[parsep=0.25ex]
	\item \zh{\textbf{实现长序列训练架构：}为 Causal-ViT 实现 Context Parallel，为 GDN 实现 Ulysses All-to-All、为 MoDA 时序注意力实现跨分片 halo 通信；支持图片--视频混训分片与损失归约，解除 CP 路径 batch size$=$1 限制，CP8 单样本训练长度从 131K 扩展至 2M visual tokens。}
	\item \zh{\textbf{支撑千卡训练稳定性：}支撑 Stage-2 千卡长视频预训练（CP8$\times$DP128、GBS 2,048、单样本最高 2,048 帧/约 2M visual tokens）；定位 CP8 短视频空分片漏发反向 All-Reduce 导致 collective 序列失配与训练 hang，以计算图锚点统一各 rank 通信序列，修复后无复现。}
	\item \zh{\textbf{MoDA 计算性能优化：}将 launch-bound、GPU 空转约 80\% 的逐帧 Python 循环重构为跨帧批处理，支持 batch size$>$1，单层/端到端提速 5.4$\times$/1.28$\times$；进一步设计 chunked packed MoDA，按 QKV、gate、attention 与 output 分阶段切分中间张量配合反向重算，在显存峰值仅小幅增加、仍支持 2M 序列训练的前提下单层再提速 3.36$\times$，并完成 BF16 单层前反向梯度对齐，真实训练 eval 表现无下降。}
	\item \zh{\textbf{分布式优化器：}定位 Muon 与 ZeRO-1/2 reduce-scatter 不兼容的根因：分区参数缺少 Newton-Schulz 所需完整归约梯度；向 DeepSpeed 上游提交切片定向传播与 owner rank 冗余更新修复，\textbf{PR\#8141 已合入 master}。}
\end{itemize}

\zh{\role{Flink 大状态计算调优}{2026.03 -- 2026.05}}
\zh{{\small\itshape 华为 · 计算产品线 | AI 工程师}\par}
\begin{itemize}[parsep=0.25ex]
	\item \zh{\textbf{状态访问链路优化：}面向 Flink 1.16.3 + RocksDB 大状态场景，构建 Value/Map 双层堆内缓存与 MapState EMPTY/SINGLE 快照；进一步以 Mailbox 攒批驱动 MultiGet 异步预取和同 key 本地预聚合，并在 MemTable/SST 前加入 whole-key Bloom miss gate，减少序列化、JNI 往返与无效读取。}
	\item \zh{\textbf{在鲲鹏平台取得显著性能收益：}在亿级数据量的 Nexmark 评测中，逐 Query 平均吞吐较 RocksDB StateBackend 在鲲鹏平台提升 50\%，并通过 CDC 数据一致性检验。}
\end{itemize}

\section{\zh{项目经历}}
\zh{\role{SQLForge：执行反馈驱动的 Text-to-SQL Agent 后训练\textnormal{\normalfont\itshape{ | QLoRA / GRPO / vLLM}}}{2026.01 -- 2026.05}}
\begin{itemize}[parsep=0.25ex]
	\item \zh{\textbf{执行反馈后训练与交互式自纠错：}构建 Agent 后训练管线——生成多 teacher ReAct 轨迹并用真实执行结果回填 Observation，QLoRA SFT 后接自研 GRPO 循环与格式/合法性/结构/部分执行/最终正确性稠密奖励；推理侧设计 Direct SQL + 执行失败后 ReAct 多轮交互，沙箱执行草稿 SQL 并按 Observation 修正，Spider 执行正确率 78.6\%，较 ReAct SFT direct-only 71.7\% 提升 6.9pp（Hard 难度 +13.4\%）。}
\end{itemize}


\section{\zh{科研经历}}
\zh{\role{DBPecker: 分布式数据库复合异常诊断系统}{PVLDB 2025 · 共同一作}}
\begin{itemize}[parsep=0.25ex]
	\item \zh{\textbf{复合异常诊断：}构建跨节点 compound anomaly graph，以相关性加权 PageRank 融合节点影响力与多标签异常概率，输出按影响排序的 $\langle$node, anomaly$\rangle$ 序列；通过结构因果模型干预识别跨 OS/数据库层根因指标，并在 OceanBase 集群故障注入完成端到端验证。}
\end{itemize}

\section{\en{Skills}\zh{技能}}
\begin{itemize}[parsep=0.25ex]
	\item \zh{\textbf{大模型系统：}Megatron-LM、DeepSpeed、Transformer Engine FP8、Triton、PyTorch、vLLM、SFT/GRPO/ReAct；千卡预训练性能、显存与 kernel 优化；}
	\item \zh{\textbf{编程与工具：}熟练使用 C++/Python 与 Claude Code、Codex；了解 Java/Go/Rust；熟悉 Linux、Git、Docker、CMake；}
	\item \zh{\textbf{英语与荣誉：}雅思 7.5（阅读/听力 8.5）、六级 594；中国机器人及人工智能大赛百度 Apollo 自动驾驶仿真赛全国一等奖。}
\end{itemize}

\end{document}
